%================================================================
\section{The estimand}\label{sec:estimand}

\subsection{Policies are assignment vectors}

The village has $N=\NAgents$ people, a fixed undirected graph with adjacency
matrix $A$, and covariates $X$. An assignment is a vector
$\dvec\in\{0,1\}^N$, where $d_i=1$ means that person $i$ is seeded. A policy is a
rule that produces an assignment. We use three:
\begin{itemize}
  \item $\pi_0$, the status quo, seeds no one: $\dvec=\dzero$.
  \item $\pi_1$, the pilot, seeds $K=\KSeeded$ people drawn uniformly at random
        without replacement, so $K/N=\SeedShare$.
  \item $\pi_2$, the policy to predict, seeds the $K$ people with the highest
        degree, with ties broken toward the lower index. Given the graph it is a
        single fixed vector, which we write $\dtwo$.
\end{itemize}

A parameter point $c$ indexes a complete data-generating process: for each $c$,
each draw $\omega$ of the simulation noise, and each assignment $\dvec$, the
kernel returns the outcomes of everyone at the horizon $T$
(Section~\ref{sec:model}). Under interference, the outcome of person $i$ depends
on the whole vector, $Y_{iT}=Y_{iT}(\dvec)$, so the effect of seeding person $i$
alone is not defined until the assignment of everyone else is named. The
well-defined object is a contrast of policies,
\begin{equation}\label{eq:theta}
  \theta(\pi,\pi';c)
  \;=\;
  \mathbb{E}_c\!\left[\bar Y_T(\pi)\right]
  \;-\;
  \mathbb{E}_c\!\left[\bar Y_T(\pi')\right],
\end{equation}
where $\bar Y_T$ is the mean outcome over people and the expectation is over the
simulation noise and, for $\pi_1$, over the random assignment. This is an overall
effect of one assignment rule against another in the sense of
\citet{halloran1995} and \citet{hudgens2008}. The code defines three such
contrasts:
\begin{align}
  \theta_1(c) &= \theta(\pi_1,\pi_0;c), \label{eq:theta1}\\
  \theta_2(c) &= \theta(\pi_2,\pi_0;c)
     = \mathbb{E}_c\bigl[\bar Y_T(\dtwo)\bigr]-\mathbb{E}_c\bigl[\bar Y_T(\dzero)\bigr],
     \label{eq:theta2}\\
  \prem(c) &= \theta_2(c)-\theta_1(c). \label{eq:premium}
\end{align}
The headline estimand is $\theta_2$, the effect of seeding the hubs against no
campaign, and $\prem$ is the premium of targeting over random seeding. In the
code each expectation is a mean over $\RepsEstimand$ paired replications: rep $r$
of every arm shares one state stream (the starting state and every shock), so the
contrasts between arms carry little Monte Carlo noise. The estimands are defined
on the terminal mean $\bar Y_T$.

Equation~\eqref{eq:theta2} makes the main point of this section. The headline
estimand compares two specific vectors, $\dtwo$ and $\dzero$. It is a property of
those two vectors, of the graph and of the kernel at $c$. Any data that never
contain $\dtwo$ can speak to it only through assumptions that connect what they
do contain to what $\dtwo$ would produce.

\subsection{What a randomized pilot reveals}

A pilot draws one assignment from $\pi_1$. Under that design every vector with
exactly $K$ ones has the same probability, so the pilot realizes the policy vector
with probability
\begin{equation}
  \Pr_{\pi_1}\!\bigl(\dvec=\dtwo\bigr)=\binom{N}{K}^{-1}
  =\binom{\NAgents}{\KSeeded}^{-1}.
\end{equation}
That probability is positive and negligible. The pilot realizes the policy with
negligible probability; in practice it never does, and positivity holds for the
policy vector only in this negligible sense. In practice, therefore, an estimate
of $\theta_2$ from the pilot needs an exposure map and transport.

What the pilot does reveal is $\theta_1$, the effect of the pilot's own design.
The difference between the mean outcome in the pilot and the mean outcome in the
status quo estimates $\theta_1$ without further assumptions. In the planted
exercise this holds exactly by construction: the pilot and status-quo targets
share their state streams, so their difference at $\cstar$ equals
$\theta_1(\cstar)$ (Section~\ref{sec:targets}). The question is whether, and
under which assumptions, one can go from the pilot to $\theta_2$.

\subsection{What an estimate from the pilot must assume}\label{sec:assumptions}

The standard route goes through exposures \citep{aronow2017,leung2020}. Let
$k_i=\max(\sum_j A_{ij},1)$ and let $s_i(\dvec)=\sum_j A_{ij}d_j/k_i$ be the
share of $i$'s neighbors who are seeded. The code uses four exposure strata,
\begin{align*}
  \Sone &= \{i: d_i=1\}, &
  \Sa &= \{i: d_i=0,\ s_i=0\}, \\
  \Sb &= \{i: d_i=0,\ 0<s_i<\ShareCut\}, &
  \Sc &= \{i: d_i=0,\ s_i\ge \ShareCut\},
\end{align*}
and we write $\Strata=\{\Sone,\Sa,\Sb,\Sc\}$ and $e_i(\dvec)\in\Strata$ for the
stratum of person $i$ under $\dvec$. Two assumptions take the pilot to
$\theta_2$.

\begin{assumption}[Exposure map]\label{ass:exposure}
The expected outcome of each person depends on the assignment only through the
person's stratum: $\mathbb{E}_c[Y_{iT}(\dvec)]=\mathbb{E}_c[Y_{iT}(\dvec')]$
whenever $e_i(\dvec)=e_i(\dvec')$.
\end{assumption}

Let $\mu_S$ be the exposure-$S$ mean outcome that the pilot measures: the
inverse-probability-weighted (H\'ajek) mean over the people who land in $S$, each
weighted by the inverse of the probability that the pilot's lottery puts them
there \citep{horvitz1952,aronow2017}.

\begin{assumption}[Transport]\label{ass:transport}
For each stratum $S$, the exposure-specific outcome $\mu_S$ measured on the
pilot's people applies to the people who would hold exposure $S$ under $\dtwo$.
\end{assumption}

Under the status quo nobody is seeded, so everyone is in $\Sa$. Every person can
land in $\Sa$ under the pilot, so under Assumption~\ref{ass:exposure} the mean
$\mu_{\Sa}$ is the status-quo mean of the whole village. Under the policy, a known
share $w_S=n_S(\dtwo)/N$ of the village holds exposure $S$; the analyst knows the
graph, so the shares are known constants that sum to one. Under both assumptions,
\begin{equation}\label{eq:transport}
  \theta_2(c)
  \;=\;\sum_{S\in\Strata} w_S\,\mu_S-\mu_{\Sa}
  \;=\;\sum_{S\neq \Sa} w_S\,\bigl(\mu_S-\mu_{\Sa}\bigr).
\end{equation}
Every term on the right can be estimated from the pilot, because the pilot's
lottery is known and the probability that each person lands in each stratum can
be computed before any outcome is seen. Equation~\eqref{eq:transport} is the
target of the analyst A1 in Section~\ref{sec:analysts}.

Neither assumption holds exactly in this village, and the paper does not pretend
that it does. The kernel's peer term uses neighbors' lagged outcomes, so the
effect of a seed travels beyond its own neighbors; a person in $\Sa$ in the pilot
can still be lifted by a seed two steps away. The exposure map therefore misses
part of the interference. Transport fails in a second way: under $\dtwo$ the
seeded people are the hubs, while under the pilot they are a random draw, and
hubs need not respond like a random draw. The analyst A2 relaxes this second
failure by splitting the village into the hubs and everyone else and transporting
only within each part; for the hubs, its term compares the hubs' own seeded and
unexposed outcomes in the pilot. Neither analyst identifies $\theta_2$ from the
pilot.
Without assumptions on how outcomes respond to the assignments of others, the
data from one randomized design bound the effects of other assignment rules only
loosely \citep{manski2013}. The analysts below differ in which assumptions they
add, and Section~\ref{sec:results} measures what each assumption costs and buys.

The analyst who assumes no interference at all (A0 below) avoids the exposure map
and transport, but only by assuming that a person's outcome depends on her own
seeding alone \citep{rubin1980}. Under that assumption, $\theta_2$ equals the
seeded share times the effect of own seeding. Under interference it does not, and
the gap between the two is measurable here (Appendix~\ref{app:gap}).
